\documentclass[11pt,a4paper]{article}
\usepackage[T1]{fontenc}
\usepackage{lmodern,amsmath,amssymb,amsthm,mathtools}
\usepackage[margin=25mm,headheight=15pt]{geometry}
\usepackage{booktabs,array,tabularx,longtable}
\usepackage[dvipsnames]{xcolor}
\usepackage{microtype,enumitem,fancyhdr,needspace}
\usepackage[colorlinks=true,linkcolor=MidnightBlue,citecolor=MidnightBlue,urlcolor=MidnightBlue]{hyperref}
\usepackage{bookmark}
\hypersetup{pdftitle={Recirculation and Diversion Accounting on a Conditional Tri-Octagon Carrier - v0.1},pdfauthor={Hilmir Frímann Halldórsson},pdfsubject={Exact source-off return and the limits of central-energy interpretation. Author-authorized repository research edition; GPT reviewed the analytical investigation, Codex prepared and checked the manuscript.}}
\setlength{\parskip}{4pt}
\setlength{\parindent}{0pt}
\setlength{\emergencystretch}{3em}
\setlist{itemsep=2pt,topsep=4pt}
\pagestyle{fancy}\fancyhf{}
\fancyhead[L]{\small Recirculation and diversion accounting}
\fancyhead[R]{\small v0.1}\fancyfoot[C]{\thepage}
\renewcommand{\headrulewidth}{0.3pt}
\numberwithin{equation}{section}
\newtheorem{theorem}{Theorem}[section]
\newtheorem{proposition}[theorem]{Proposition}
\newtheorem{corollary}[theorem]{Corollary}
\theoremstyle{remark}\newtheorem{remark}[theorem]{Remark}
\newcommand{\R}{\mathbb R}
\newcommand{\dd}{\,\mathrm d}
\newcommand{\eps}{\varepsilon}
\newcommand{\norm}[1]{\left\|#1\right\|}
\newcommand{\re}{\operatorname{Re}}
\newcommand{\im}{\operatorname{Im}}
\newcommand{\ez}{e_z}
\newcolumntype{Y}{>{\raggedright\arraybackslash}X}
\newcolumntype{P}[1]{>{\raggedright\arraybackslash}p{#1}}
\begin{document}
\hypersetup{pageanchor=false}
\begin{titlepage}
\vspace*{6mm}
{\small\color{MidnightBlue} TRI-OCTAGON MATHEMATICAL RESEARCH\par}
\vspace{8mm}
{\LARGE\bfseries\raggedright Recirculation and Diversion Accounting\\
on a Conditional Tri-Octagon Carrier\par}
\vspace{6mm}
{\large Exact source-off return and the limits\\ of central-energy interpretation\par}
\vspace{7mm}
{\large Hilmir Fr\'imann Halld\'orsson\par}
\vspace{3mm}
{Author-authorized repository research edition v0.1\quad 3 November 2026\par}
\vspace{5mm}
\begin{abstract}
A separately adopted periodic one-dimensional carrier associated with the Tri-Octagon model admits exact source-off recirculation. Its normalized characteristic profiles travel in opposite directions, preserve their separate quadratic storage, and return to the same periodic coordinates after one circuit time. At a fixed section, one full source-off circuit carries each directional storage past the cut once. Gross throughput equals total storage, while net throughput is the difference of the directional stocks; repeated circuits recount circulating storage. Complementary-arc balances distinguish local departure from total loss, and a same-drive preparation comparison retains the signed nature of source work. We then separate a hypothetical diverted share from a finite-window return deficit: the same shortfall is compatible with diversion or temporary route storage. An unloaded exterior observation adds no supplying-energy debit. These results establish conditional carrier accounting, not an energetic transfer from or return to the native center. Central-route registration, receiving states and reciprocal native exchange remain unassigned, and no numerical diversion fraction is inferred.
\end{abstract}
\vfill
{\small Research author and questions: Hilmir Fr\'imann Halld\'orsson (pzychozen).
GPT reviewed and accepted the underlying analytical investigation of 2 November 2026, as recorded in the publication authorization.
Codex prepared this manuscript and checked its fidelity, build and final pages.
GPT has not separately inspected this final typeset manuscript.
This is repository research, without external peer review, journal acceptance, DOI or experimental validation.}
\end{titlepage}
\setcounter{page}{2}\hypersetup{pageanchor=true}
\section{Question, scope and the inherited carrier}\label{sec:scope}
The author's intended picture is circulation: a central starting or meeting region, paired upward and downward cones, outward movement curving around an intended tight toroid, and inward return. A possible portion might go elsewhere instead of returning. The suggested \emph{15--20\%, possibly more} is an unverified intuition, not a measurement, fitted coefficient or accepted bound. Flower and heart-field analogies communicate the picture; they supply neither field lines nor a physiological or electromagnetic mechanism.

The accepted investigation \cite{assessment} distinguishes what already recirculates mathematically from the missing transfer to this intended central route. This note presents its written deductions as a companion to the native intensity and crossings paper \cite{intensity}. It does not choose a destination, loading law, energy calibration or physical time bridge. In particular, leaving a local arc and leaving the whole carrier are different events.

\subsection{Conditional state and source profiles}\label{sec:carrier}
We reuse C1 from the paired-driving paper \cite[carrier and foundations sections]{pairedII}, writing $r_c,g_c,L_c,\eps_c$ for that source's $r,g,L,\eps$ to avoid collisions with native and geometric notation. On $s\in\R/L_c\mathbb Z$, with $c>0$, the adopted equations are
\begin{equation}\label{eq:C1}
 p_t+c(r_c)_s=g_c,\qquad (r_c)_t+cp_s=0,\qquad
 g_c(s,t)=b_+(s)f_+(t)+b_-(s)f_-(t).
\end{equation}
The real preparation $(p_0,r_{c,0})$ is independently specified. At the registered contacts $s_\pm$, the profiles are translates of a periodic tent:
\begin{equation}\label{eq:tent}
 b_\eps(s)=\eps_c^{-1}\left(1-\frac{|\operatorname{wrap}_{L_c}(s)|}{\eps_c}\right)_+,
 \qquad b_\pm(s)=b_\eps(s-s_\pm),\qquad 0<\eps_c<\eta.
\end{equation}
Here the shorter contact separation is $2\eta$. The profiles have disjoint supports, unit integral, and belong to $H^1_{\rm per}$. The width is arclength width, not a tube radius. The $f_\pm$ are distributed source strengths, not prescribed values of $p$ at endpoints. The circle has no endpoints.

The inherited well-posedness theorem gives, for periodic $H^1$ preparation and continuous drives on finite intervals, a unique solution in
\begin{equation}\label{eq:regularity}
 C([0,T_*];H^1_{\rm per})^2\cap C^1([0,T_*];L^2)^2.
\end{equation}
For $L^2$ preparation and $L^1_tL^2_s$ forcing the characteristic formula defines a unique $C_tL^2_s$ mild solution. The source's work theorem specifically extends its integrated global identity, with derivative almost everywhere, to finite-profile time-$L^2$ drives. These include the admitted finite sample/hold records. We inherit these results rather than solving a new history.

\subsection{Storage and signed work}\label{sec:workaccount}
The definitions and balances are
\begin{equation}\label{eq:account}
\begin{gathered}
 e=\tfrac12(p^2+r_c^2),\qquad j=cpr_c,\qquad
 E(t)=\int_0^{L_c}e\dd s,\\
 y_\pm(t)=\int_0^{L_c}b_\pm p\dd s.
\end{gathered}
\end{equation}
\begin{equation}\label{eq:balance}
 e_t+j_s=pg_c,\qquad E'(t)=f_+y_++f_-y_-,
\end{equation}
\begin{equation}\label{eq:integrated}
 E(b)-E(a)=\int_a^b\bigl(f_+y_++f_-y_-\bigr)\dd t.
\end{equation}
The local identity is distributional at the strong regularity above; the global derivative is classical under those data and almost everywhere in the stated mild extension. The integrated form is the operative weak global account. Source-off norm conservation and profile recurrence below also hold at $L^2$ preparation.

\textbf{Trace convention.} All fixed point-cut statements in this note use spatial $H^1$ regularity, with the associated solution/time regularity just stated. In one dimension this supplies continuous spatial representatives and the weak product rule. We make no unrestricted point-flux claim for arbitrary $L^2$ states.

\subsection{The existing native input interface}\label{sec:input}
C1 is not the native recurrence. For one named readout family $v_{\rm in}=M$ or $C$, the source nevertheless already defines a conditional chain
\begin{equation}\label{eq:input}
 \Omega,\text{ clock/settings}\ \longmapsto\ v_{\rm in}
 \ \longmapsto\ A_{\rm reg}v_{\rm in}
 \ \longmapsto\ f_\pm=\kappa_{\rm in}n_\pm^TA_{\rm reg}v_{\rm in},
\end{equation}
where $A_{\rm reg}=\lambda_gR_z(\delta_g)$ and $\kappa_{\rm in}>0$ \cite[typed input chain]{pairedII}. The paired projection loses the component along $v=e_z\times u$. Its positive gain converts data units to source strength; it is not energy efficiency. A finite record can be held over common intervals $[j\Delta t,(j+1)\Delta t)$, with symbolic $\Delta t>0$ and an end gate. That is an interface assumption, not intrinsic native continuous time.

The missing relationship is therefore not every input map. It is reciprocal energetic exchange: \eqref{eq:C1} has no paired native debit for positive source work, native credit for negative work, or receiving-load backreaction. The native index, readout clock, carrier time and geometric parameter remain distinct.

\section{Exact source-off return and section throughput}\label{sec:return}
For $\sigma\in\{+1,-1\}$, define normalized characteristic fields, their densities and stocks by
\begin{equation}\label{eq:chars}
 w_\sigma=\frac{p+\sigma r_c}{\sqrt2},\qquad
 \rho_\sigma=\frac{w_\sigma^2}{2},\qquad
 E_\sigma(t)=\int_0^{L_c}\rho_\sigma\dd s.
\end{equation}
Adding and subtracting \eqref{eq:C1} gives the inherited transport equation. Multiplication by $w_\sigma$ gives its directional account:
\begin{align}
 (\partial_t+\sigma c\partial_s)w_\sigma&=g_c/\sqrt2,\label{eq:transport}\\
 \partial_t\rho_\sigma+\sigma c\partial_s\rho_\sigma
 &=w_\sigma g_c/\sqrt2.\label{eq:directional}
\end{align}
The product calculation has the distributional meaning specified in Section~\ref{sec:workaccount}. Inversion of the orthogonal transform gives
\begin{equation}\label{eq:netgross}
 e=\rho_++\rho_-,\qquad
 j=c(\rho_+-\rho_-),\qquad
 j_{\rm gross}:=c(\rho_++\rho_-)=ce,\qquad |j|\le j_{\rm gross}.
\end{equation}
The signed directional fluxes are $+c\rho_+$ and $-c\rho_-$. Their nonnegative magnitudes add to the gross flux, which is generally not $|j|$.

\begin{theorem}[Conditional source-off recurrence]\label{thm:return}
On an entirely source-off interval, $E_+$ and $E_-$ are separately conserved. If $[t,t+T_{\rm loop}]$ is entirely source-off, with
\begin{equation}\label{eq:period}
 T_{\rm loop}=L_c/c,
\end{equation}
then $(p,r_c)(\cdot,t+T_{\rm loop})=(p,r_c)(\cdot,t)$. The equality holds in $L^2$ for mild preparation and in the continuous spatial representatives at $H^1$ regularity. $T_{\rm loop}$ is a period, not necessarily the smallest one.
\end{theorem}
\begin{proof}
The inherited characteristic representation \cite[well-posedness theorem]{pairedII} specializes, on a source-off interval starting at $a$, to
\begin{equation}\label{eq:translation}
 w_\sigma(s,t)=w_\sigma(s-\sigma c(t-a),a).
\end{equation}
Thus the $+$ profile travels toward increasing $s$ and the $-$ profile toward decreasing $s$, both at speed $c$. Periodic translation is an $L^2$ isometry, so it preserves each squared norm and hence each $E_\sigma$. Translation through $-\sigma cT_{\rm loop}=-\sigma L_c$ is the identity on the periodic domain. Inverting the characteristic transform proves recurrence of $(p,r_c)$. No newly evaluated wave solution is needed.
\end{proof}

\begin{theorem}[One-circuit section account]\label{thm:throughput}
At a declared fixed cut $s_*$, under the trace convention in Section~\ref{sec:workaccount}, define
\begin{equation}\label{eq:Q}
 Q_\sigma(s_*;a,b)=\int_a^b c\rho_\sigma(s_*,u)\dd u.
\end{equation}
If the entire interval $[t,t+T_{\rm loop}]$ is source-off, then
\begin{align}
 Q_\sigma(s_*;t,t+T_{\rm loop})&=E_\sigma(t),\label{eq:Qdirection}\\
 Q_{\rm net}:=\int_t^{t+T_{\rm loop}}j(s_*,u)\dd u
 &=E_+(t)-E_-(t),\label{eq:Qnet}\\
 Q_{\rm gross}:=\int_t^{t+T_{\rm loop}}j_{\rm gross}(s_*,u)\dd u
 &=E_+(t)+E_-(t)=E(t).\label{eq:Qgross}
\end{align}
\end{theorem}
\begin{proof}
Apply \eqref{eq:translation} with starting time $t$ in \eqref{eq:Q}. The change of variables $x=s_*-\sigma c(u-t)$ covers the circle once in time $L_c/c$. For $\sigma=+1$ the reversed limits cancel the minus sign from $\dd x=-c\dd u$; for $\sigma=-1$ the orientation is increasing. In either case periodicity gives
\[
 \int_t^{t+L_c/c}c\rho_\sigma(s_*-\sigma c(u-t),t)\dd u
 =\int_0^{L_c}\rho_\sigma(x,t)\dd x.
\]
This proves \eqref{eq:Qdirection}. The difference and sum, using \eqref{eq:netgross}, give \eqref{eq:Qnet} and \eqref{eq:Qgross}.
\end{proof}

The theorem distinguishes counterpropagation from net transport. Equal positive directional stocks have zero full-circuit net cut amount and positive gross throughput. It does not require a preferred direction or nonzero flux for every preparation.

Over $m$ whole source-off circuits, $m$ a positive integer, adding \eqref{eq:Qgross} gives gross throughput $mE$, while storage remains $E$. This recounts circulating storage, not $mE$ of new emission. A statement that all returns requires a declared profile or same-direction section-turnover meaning and a positive denominator. Neither theorem identifies tagged particles, a launch cohort in a superposed driven field, a native-state recurrence, or receipt at the spatial center.

\begin{corollary}[Zero source-off total-storage depletion]\label{cor:depletion}
If $E(t)>0$ and $[t,t+T_{\rm loop}]$ is source-off, then
\begin{equation}\label{eq:depletion}
 d_E=\frac{E(t)-E(t+T_{\rm loop})}{E(t)}=0.
\end{equation}
At $E(t)=0$ this ratio is undefined, while conservation still holds.
\end{corollary}
\begin{proof}
The sum $E=E_++E_-$ is conserved. A proposed law
$E(t+T_{\rm loop})=(1-d)E(t)$ with $d>0$ would therefore contradict conservation when $E(t)>0$.
\end{proof}
The same storage obstruction holds on every source-off interval, not only a full circuit. It is a theorem about unmodified C1, whose inspected equations contain no intrinsic sink, damping or leakage branch. It is not a physical measurement showing that the author's intended circulation is lossless. The source's statement that autonomous circulation has not been established concerns the missing physical and reciprocal core interpretation \cite[conclusion]{pairedII}; it does not negate homogeneous mathematical profile recurrence.

\section{Driven work and internal arc transfers}\label{sec:driven}
\subsection{Preparation matters even at fixed initial storage}
Let
\begin{equation}\label{eq:Wsrc}
 W_{\rm src}[a,b]=\int_a^b\bigl(f_+y_++f_-y_-\bigr)\dd t.
\end{equation}
By \eqref{eq:integrated}, positive net work increases total storage and negative net work reduces it. Zero net work can conceal contributions of both signs. This account does not identify cut flow with source work.

Equal drives and equal full preparation give the same state by uniqueness. Drives alone do not determine its storage history, even when the initial total storage is fixed.
\begin{proposition}[Same drives, equal initial stock, different work signs]\label{prop:preparation}
For continuous initial forcing $g_0=g_c(\cdot,a)\ne0$ and $A>0$ with time dimension, the admissible preparations
\begin{equation}\label{eq:preparations}
 p_0^{(\pm)}=\pm A g_0,\qquad r_{c,0}^{(\pm)}=0
\end{equation}
have equal initial total storage and opposite initial derivatives under the same drives:
\begin{equation}\label{eq:preparationwork}
 E^{(\pm)}(a)=\tfrac12A^2\norm{g_0}_2^2,\qquad
 (E^{(\pm)})'(a)=\pm A\norm{g_0}_2^2.
\end{equation}
\end{proposition}
\begin{proof}
The finite tent profiles make $g_0$ an $H^1$ function, so both preparations satisfy the inherited hypotheses. Substitution in the storage definition gives the same squared norm. Substitution in $E'(a)=\int p_0g_0\dd s$ gives the stated signs. Distinct initial derivatives imply distinct storage histories.
\end{proof}
This is a symbolic comparison of permitted preparation data, not a newly computed trajectory. It uses the source's signed-work example \cite[work theorem and preparation example]{pairedII}. The nonzero-forcing assumption matters: source-off states with equal initial storage have the same constant total-storage history, though their fields may differ.

Negative work removes C1 storage through the existing drive interface. It does not name an exterior receiver or credit the native core. Conversely positive work has no reciprocal native debit. Source strengths squared do not replace $f_\pm y_\pm$. Moreover, linearity of the field equation does not imply additivity of quadratic storage: superposed responses retain cross terms \cite[quadratic storage account]{pairedII}. A launch amount attributable to a new drive in a pre-existing field therefore requires explicit bookkeeping.

\subsection{Two host cuts and their complement}\label{sec:arcs}
The accepted registration gives, with its outward host-plane signs \cite[conditional signed cut powers]{intensity},
\begin{equation}\label{eq:cuts}
 P_{\rm out,+}=-j(s_+,t),\qquad
 P_{\rm out,-}=+j(s_-,t).
\end{equation}
Let $I=[s_+,s_-]$ be the short arc and $K=[s_-,s_++L_c]$ its periodic complement. Put
\begin{equation}\label{eq:arcterms}
 B=j(s_-,t)-j(s_+,t),\qquad
 E_J=\int_J e\dd s,\qquad \dot W_J=\int_J pg_c\dd s
 \quad(J=I,K).
\end{equation}
Here $B=P_{\rm out,+}+P_{\rm out,-}$ is a signed internal boundary transfer.

\Needspace{12\baselineskip}
\begin{proposition}[Complementary-arc cancellation]\label{prop:arcs}
At the stated trace regularity,
\begin{equation}\label{eq:arcbalance}
 E_I'=\dot W_I-B,\qquad E_K'=\dot W_K+B
\end{equation}
almost everywhere. With $W_J[a,b]=\int_a^b\dot W_J\dd t$, the integrated identities are
\begin{align}
 E_I(b)-E_I(a)&=W_I[a,b]-\int_a^bB\dd t,\label{eq:arcI}\\
 E_K(b)-E_K(a)&=W_K[a,b]+\int_a^bB\dd t.\label{eq:arcK}
\end{align}
Their sum is the whole-circle work account.
\end{proposition}
\begin{proof}
The short-arc identity is inherited from the crossings paper. Integrating the local balance over the complement gives
$E_K'=\dot W_K+j(s_-)-j(s_++L_c)$. Periodicity replaces the last value by $j(s_+)$ and proves its sign in \eqref{eq:arcbalance}. Integrating in time gives \eqref{eq:arcI}--\eqref{eq:arcK}; summing cancels $B$.
\end{proof}
Source-off net loss from $I$ equals gain on $K$: storage has left a subarc and stayed on the same circle. During driving each arc also has its own work, and total loss is exactly negative net total source work. Positive $B$ alone does not establish that loss. The tents straddle the cuts, so arc-restricted work cannot generally be replaced by $W_{\rm src}$.

Positive local $j$ is incoming at the upper cut and outgoing at the lower; negative local $j$ reverses the corresponding sign. The two values need not agree at a given time. Neither geometric label is a permanent outlet. The cuts are not absorbing ends and are not automatically the boundary of a native-center control volume.

\section{What a diversion fraction would require}\label{sec:diversion}
\subsection{Stocks, powers, amounts and records}\label{sec:types}
The source's abstract dimensions are $\mathrm U,\mathrm L,\mathrm T$, with $[p]=[r_c]=\mathrm U$, $[c]=\mathrm L/\mathrm T$, $[b_\pm]=\mathrm L^{-1}$ and $[f_\pm]=\mathrm U\mathrm L/\mathrm T$ \cite[dimension table]{pairedII}. No SI unit or physical time calibration is selected. Table~\ref{tab:types} distinguishes existing quantities from proposed route accounts.

{\small
\begin{longtable}{@{}P{.24\textwidth}P{.29\textwidth}P{.41\textwidth}@{}}
\caption{Existing quantities and denominator limitations.}\label{tab:types}\\
\toprule Quantity & Type or dimension & Interpretation limit\\\midrule
\endfirsthead
\toprule Quantity & Type or dimension & Interpretation limit\\\midrule
\endhead
\bottomrule\endfoot
$E(t),E_\pm(t)$ & Nonnegative stocks; $\mathrm U^2\mathrm L$ & Not automatically native launch energy. The depletion ratio requires $E>0$.\\[4pt]
$\int_a^b j(s_*,t)\dd t$ & Signed net cut amount; $\mathrm U^2\mathrm L$ & Can cancel to zero or be negative; not a default positive fraction denominator.\\[4pt]
$Q_\pm(s_*;a,b)$ & Nonnegative directional amounts; $\mathrm U^2\mathrm L$ & Require a cut, direction, window and positive reference amount for normalization.\\[4pt]
$\int_a^b j_{\rm gross}(s_*,t)\dd t$ & Nonnegative gross amount & Counts both directions and repeats storage on later circuits.\\[4pt]
$W_{\rm src}[a,b]$ & Signed work; $\mathrm U^2\mathrm L$ & Can remove rather than supply storage. No recipient is assigned.\\[4pt]
$j;\ \int_0^{L_c}j\dd s$ & $\mathrm U^2\mathrm L/\mathrm T$; $\mathrm U^2\mathrm L^2/\mathrm T$ & The spatial integral has an extra length dimension: it is neither a section power nor an energy amount.\\[4pt]
Native $I,V$, decoded vectors & Native intensity, potential or based-vector data & No canonical conversion into a delivered carrier-energy amount.\\[4pt]
C6 $Y_X$ & Delayed amplitude record; $\mathrm U$ & No receiving-energy account or power meaning for its square.\\[4pt]
$Q_{\rm launch},Q_{\rm return},Q_{\rm div}$ & Proposed nonnegative transfers & Intended route, destination, denominator and event accounting remain unassigned.\\
\end{longtable}}

The ratio $Q_{\rm div}/Q_{\rm launch}$ would describe a per-pass diverted share. In contrast, $1-Q_{\rm return}/Q_{\rm launch}$ would describe a finite-window return deficit. Neither definition is adopted here. Eventual non-return is a further limiting question, not a synonym for either finite-window ratio.

\subsection{A hypothetical route ledger}\label{sec:ledger}
Suppose only for accounting that a future route region and interval $[a,b]$ have been specified. Let the $Q$ terms below be nonnegative transfers with declared orientations, let $W_{\rm route}$ be signed non-boundary work, and let $\Delta S_{\rm route}$ be the signed route-inventory change. A compatible balance must read
\begin{equation}\label{eq:route}
\begin{split}
 Q_{\rm launch}+Q_{\rm other,in}+W_{\rm route}
 ={}&Q_{\rm return}+Q_{\rm div}\\
 &+Q_{\rm other,out}+\Delta S_{\rm route}.
\end{split}
\end{equation}
This is hypothetical necessary bookkeeping, not a new dynamics law. Terms need common units and windows and a rule for classifying repeated crossing events. A distributed source already booked as $W_{\rm route}$ cannot also be counted as boundary launch work without changing the accounting partition.

With no other transfers and $Q_{\rm launch}>0$, a proposed $f_{\rm div}=Q_{\rm div}/Q_{\rm launch}$ satisfies, by rearrangement,
\begin{equation}\label{eq:returnratio}
 \frac{Q_{\rm return}}{Q_{\rm launch}}
 =1-f_{\rm div}
  +\frac{W_{\rm route}-\Delta S_{\rm route}}{Q_{\rm launch}}.
\end{equation}
The simplified relation $Q_{\rm return}=(1-f_{\rm div})Q_{\rm launch}$ therefore requires $W_{\rm route}=\Delta S_{\rm route}$, for example both zero, together with no other transfers and compatible event accounting. Under that restriction the nonnegative amounts imply $0\le f_{\rm div}\le1$. With other work or inventory release a raw ratio need not obey this upper bound, and would not be a simple share of newly launched energy.

\begin{proposition}[A return deficit does not identify diversion]\label{prop:deficit}
For $Q>0$ and $0<d<1$, a finite-window launch $Q$ and return $(1-d)Q$, with zero other transfers and work, satisfy \eqref{eq:route} under either assignment
\begin{align}
 Q_{\rm div}=dQ,\quad &\Delta S_{\rm route}=0,\label{eq:divwitness}\\
 Q_{\rm div}=0,\quad &\Delta S_{\rm route}=dQ.\label{eq:storewitness}
\end{align}
Thus this return deficit alone does not distinguish diversion from temporary storage or delay.
\end{proposition}
\begin{proof}
In each assignment, the sum of return, diversion and inventory change is $(1-d)Q+dQ=Q$, equal to launch.
\end{proof}
These are account-level witnesses, not two constructed native or carrier trajectories. To test a diverted share requires an actual destination state or receiving-power account. Attribution of supply to the native core additionally needs a commensurate native storage candidate, a time bridge and reciprocal debit/credit or backreaction. Neither native intensity $I$ nor potential $V$ is chosen as that storage. Preparation, baseline and an attribution convention are needed where existing waves and quadratic interference are present.

\subsection{Internal destinations and the total boundary}\label{sec:internal}
If the destination were internal to an enlarged core--route--destination union, the compatible hypothetical special case would be
\begin{align}
 \Delta S_{\rm route}&=Q_{\rm launch}-Q_{\rm return}-Q_{\rm div}
                         +W_{\rm route},\label{eq:internalroute}\\
 \Delta S_{\rm core}&=-Q_{\rm launch}+Q_{\rm return}+W_{\rm core},
                         \label{eq:internalcore}\\
 \Delta S_{\rm dest}&=Q_{\rm div}+W_{\rm dest}.\label{eq:internaldest}
\end{align}
Addition cancels all three internal transfers:
\begin{equation}\label{eq:union}
 \Delta(S_{\rm route}+S_{\rm core}+S_{\rm dest})
 =W_{\rm route}+W_{\rm core}+W_{\rm dest}.
\end{equation}
Diversion from a return path therefore need not be loss from this enlarged total system. If the destination lies outside the chosen boundary, its transfer instead belongs to that total's outgoing account. A dissipative interpretation needs its own energy/heat account if total conservation is claimed. None of these identities selects the author's boundary or installs receiving dynamics.

\subsection{Per-pass removal is not eventual efficiency}\label{sec:compounding}
Under the additional hypothetical assumptions of a constant per-pass removed share $f$, compatible work/inventory accounting, and no re-entry of removed energy,
\begin{equation}\label{eq:perpass}
 Q_{m+1}=(1-f)Q_m,\qquad Q_m=(1-f)^mQ_0.
\end{equation}
Repeated substitution gives the second formula. Telescoping the removals gives
\begin{equation}\label{eq:cumulative}
 Q_{\rm removed,cumulative}(m)=Q_0[1-(1-f)^m].
\end{equation}
For $0<f<1$ the remaining amount tends to zero as $m\to\infty$. Thus an assumed per-pass share cannot be silently reported as a one-time or eventual efficiency. Re-entry, changing shares, forcing and finite observation windows change the account.

No accepted theorem confines the intended route share to $[0.15,0.20]$, establishes a value above it, or establishes any actual nonzero leak. The source-off C1 depletion ratio is zero for positive storage; the central-route diversion ratio is unidentifiable from the current definitions. A future expanded or intermittently diverted system is not refuted by this scope-limited theorem. The author's range could only be a future hypothesis or, after defining and measuring both amounts, an empirical ratio. No value is evaluated or fitted here.

\section{Exterior observation and central-route interpretation}\label{sec:interpretation}
\subsection{The unloaded C6 history map}\label{sec:C6}
For a fixed off-circle point $X$, an orthonormal observation frame $B_{\rm obs}$, distance $d_X(s)=\norm{X-\Gamma(s)}$, independent speed $c_e>0$ and dimensionless gain $g_o>0$, the adopted C6 rule is \cite[exterior section]{pairedII}
\begin{align}
 U(X,t)&=\frac{g_o}{4\pi}\int_0^{L_c}
       \frac{p(s,t-d_X(s)/c_e)}{d_X(s)}\dd s,\label{eq:C6scalar}\\
 V_{\rm ext}(X,t)&=\frac{g_o}{4\pi}\int_0^{L_c}
       \frac{r_c(s,t-d_X(s)/c_e)\hat t(s)}{d_X(s)}\dd s,\qquad
 Y_X=(U,B_{\rm obs}^TV_{\rm ext}).\label{eq:C6vector}
\end{align}
The source provides a bounded causal history map on its specified $L^2$ history domain. The delay, distance weight, tangent encoding and absence of loading are assumptions. The vector $r_c\hat t$ is not the carrier flux $j$, native chirality or a material current.

\begin{proposition}[Observation does not impose a debit]\label{prop:C6}
For a fixed admissible carrier history, rescaling $g_o$ rescales any nonzero C6 record without changing C1 storage or its work balance. Thus a nonzero exterior record does not, within the inspected equations, require energetic loss from C1.
\end{proposition}
\begin{proof}
Neither $g_o$, $X$ nor $Y_X$ enters \eqref{eq:C1} or \eqref{eq:balance}. Fixing preparation and drives fixes the carrier history by uniqueness. Equations \eqref{eq:C6scalar}--\eqref{eq:C6vector} depend linearly on $g_o$, so changing that gain changes the record without changing the supplying state. Adding observations also inserts no term in its balance.
\end{proof}
A delayed exterior record can therefore be nonzero during a source-off interval without violating carrier conservation. No new history is evaluated for this argument. A receiving-energy interpretation needs a source/receiver power account and a load or reciprocal reaction; the source also flags the missing ambient energy and self-force construction. Neither a radiating-looking display nor $|U|^2+\norm{V_{\rm ext}}^2$ is thereby emitted power or Poynting flux.

\subsection{What the registered geometry does and does not say}\label{sec:geometry}
The registered C1 centerline is the adopted vertical offset circle
\begin{equation}\label{eq:circle}
 \Gamma(s)=O+K_cu+R_c(\cos\chi\,u+\sin\chi\,e_z),\qquad
 \chi=s/R_c,\qquad L_c=2\pi R_c.
\end{equation}
The fixed geometry gives $K_c-R_c>0$ \cite[normal-crossing circle]{pairedI}. Hence
$(\Gamma(s)-O)\cdot u\ge K_c-R_c>0$, and the exact point O is excluded. The circle's axis is perpendicular to the central shell axis. It is a one-dimensional centerline, without an assigned tube cross section or surrounding field.

\textbf{Point versus region.} Exclusion of O alone does not prove that the circle misses every possible finite central starting or meeting region. No receiving/control region has been selected, and no registration to the intended central return route has been established. A circle alone neither specifies nor rules out a surrounding three-dimensional field. This interpretation qualification, recorded with acceptance of the underlying investigation, does not change the original proof.

The finite host-notch measurement triangles $N_k^\pm$ are different from the central broken boundaries $G_k^\pm$. The accepted dictionary is
\begin{equation}\label{eq:patchdictionary}
 \Delta_j^\sigma-O=N_{(j+1)\bmod3}^\sigma.
\end{equation}
C1 selects geometric index $j=0$, hence $N_1^+,N_1^-$. Its centerline crosses that pair normally and meets none of the other four finite patches. The source's three geometric candidates do not instantiate three carriers or a six-port coupling. Their patch normals justify \eqref{eq:cuts} \cite[geometry and conditional cuts]{intensity}.

By contrast, the central definitions give the joined segments
\begin{equation}\label{eq:G}
 G_k^+=U_k\,--\,T_{2k+1}\,--\,U_{k+1},\qquad
 G_k^-=L_k\,--\,B_{2k+1}\,--\,L_{k+1},\quad k=0,1,2.
\end{equation}
Here $U_k,L_k$ are the upper/lower fan tips and $T_i,B_i$ the corresponding ring vertices; the source indices are cyclic. These are two-segment boundaries, not flow directions or filled apertures. No closing chord, flux normal or central control volume is added. The accepted fixed-gap obstruction rules out placing an entire such arc in a host patch by a centered rotation, but does not forbid a future transfer registration \cite{fixedgap,intensity}.

The original close sphere, of radius $\sqrt{13/12-\sqrt2/2}$ about O, is a geometric reference, not a wall. It remains distinct from the extended-line sphere and the author's intended tight toroid and paired cones. The positioned $180^\circ$ support lines remain stationary. Static geometric $\tau$, native common phase, native iteration, carrier time and observer clock are different variables. Two co-present arrangements are not silently replaced by one rapidly alternating arrangement.

\subsection{Native quantities as observations, not transfer laws}\label{sec:native}
The predecessor \cite{intensity} supplies the exact ideal finite-step native intensity account. With $s_i=|\Omega_i|^2$, $D_{ij}=|\Omega_i-\Omega_j|^2$ and prestage increment
$b_i=\eps(k_i-s_i)\Omega_i+g\sum_{j\ne i}(\Omega_j-\Omega_i)$, it is
\begin{equation}\label{eq:native}
 \Delta I=2\eps\sum_i s_i(k_i-s_i)-2g\sum_{i<j}D_{ij}+\norm b^2.
\end{equation}
The mismatch and complete finite-step square remain. The selected antisymmetric exchange $J^I_{i\to j}=g(s_i-s_j)$ is channel bookkeeping, not a uniquely selected spatial current. This ideal identity is algebraic at general real coefficients; sign claims need their stated coefficient assumptions. The native potential $V$ and its fixed-domain theorem remain those of the predecessor and its cited sources; neither $V$ nor $I$ is identified with C1 storage or its supply. No historical Z trajectory or scientific check is repeated.

The accepted decoder gives based vectors
\begin{equation}\label{eq:decoder}
 v_i=\re(\Omega_i)t_i+\im(\Omega_i)e_z,\qquad (c_i,v_i),\qquad
 \norm{v_i}=\sqrt{s_i}.
\end{equation}
For the adopted frames, $n_i=-g_{i-1}$ and
$c_i+\R n_i=O+\R g_{i-1}$, where the $g_k$ here are geometric support directions. Pure common phase rotates each vector about its own stationary support line, not as a general global rigid-body pose. Full F also changes amplitudes and relative phases. Matched tangent-plane transport compares vectors; it is not motion along seams. Vector length is not travel distance \cite{phase,intensity}.

On the accepted coherent domain, the five retained coordinates
$y=(s_1,s_2,s_3,A_{12},A_{13})$ determine coherence; the accepted reconstruction additionally needs initial common phase $u_0$ for actual phase-sensitive vectors. A full recorded complex state already contains that phase. Frozen-setting staged Z loses common phase, so a genuinely common-phase-dependent outlet quantity cannot be inferred from that staged history alone. Recovering the vectors does not supply an outlet-transfer rule. Raw channel chirality $C=\re\Omega\times\im\Omega$ and its ambient attachment $T_fC$ remain distinct; clock/settings and EMA memory have their separate information roles. None is advanced in this manuscript.

\subsection{Dependency and missing exchange}\label{sec:interfaces}
Table~\ref{tab:interfaces} retains the positive input and carrier results while isolating the missing links. Recurrence of C1's state is not a feedback term into F. The central-boundary row records the requested interpretation; it does not prescribe that every future route must traverse G. No six-opening allocation law is selected from three native channels or two C1 contacts.

\begin{table}[htbp]\centering\small
\begin{tabularx}{\textwidth}{@{}P{.39\textwidth}Y@{}}\toprule
Dependency & Status\\\midrule
Native state $\to$ named readout & Defined with its family and settings; retain phase if needed.\\[3pt]
Readout $\to$ two drive strengths & Conditional registration, gain and sample/hold; not an energy calibration.\\[3pt]
Drives and preparation $\to$ C1 & Conditional well-posed state, $E$ and signed $j$.\\[3pt]
C1 $\to$ same periodic coordinates & Exact source-off profile recurrence after $L_c/c$.\\[3pt]
C1 $\to$ host cuts / C6 records & Existing two cut signs / unloaded delayed observations.\\[3pt]
Host cuts $\to$ central $G$ / intended route & Transfer registration and route assignment absent.\\[3pt]
Route $\to$ return/diversion receiver & Receiving state, oriented transfer and denominator absent.\\[3pt]
Supply/return $\to$ native debit/credit & Reciprocal balance, load and native time/unit bridge absent.\\\bottomrule
\end{tabularx}
\caption{Text dependency map; no geometric flow construction is implied.}\label{tab:interfaces}
\end{table}

\section{Conclusions and reproducibility}\label{sec:conclusion}
\textbf{Inherited native results.}
Exact intensity accounting, the separately scoped native potential results, phase reconstruction, based-vector decoding and stationary supports remain inherited. They supply no destination for a decrease or energetic receipt at a spatial center.

\textbf{Conditional C1 facts.}
The inherited C1 results provide existence, uniqueness, a characteristic representation and signed work. Its sources are distributed over the periodic domain; the registered cuts are internal sections. C6 supplies unloaded observations without a receiving balance.

\textbf{Deductions established in this investigation.}
Source-off directional stocks are conserved and characteristic profiles recur after $L_c/c$. One full-circuit section throughput equals each directional stock; gross throughput equals $E$, and repeated circuits recount it. Complementary arcs exchange storage with cancelling internal signs. Equal initial stock and equal drives can coexist with different work histories through different preparation. Hypothetical route ledgers show why return deficit does not identify diversion and why per-pass removal is not an eventual efficiency.

\textbf{Unassigned relationships.}
An energetic transfer from and return to the native center, receiving/diverting states, compatible units/time, reciprocal debit/credit, and registration to the intended central route remain unspecified. For positive storage, unmodified source-off C1 has zero total-storage depletion. The intended central-route diversion share has no supplied numerator/denominator account; neither a nonzero exterior record nor profile recurrence supplies it.

\textbf{Author's intuition.}
The central meeting region, cones, outward curved path and inward return remain the motivating picture. The possible 15--20\% diversion, or more, is neither inferred nor fitted, and is not judged true or false as a physical measurement by this conditional analysis.

The prospective modeling checkpoint is whether the destination lies inside or outside a chosen total-system boundary. That choice would change the total balance, but would still require actual receiving/transfer and reciprocal-native definitions where claimed. Neither the boundary nor those laws are chosen by this publication.

\subsection*{Evidence, attribution and reproducibility}
The underlying investigation was accepted at its commissioned scope on 2 November 2026 as recorded by the author-supplied publication order. Its original proof, ledger and order are retained unchanged in this edition's evidence directory. A separately supplied GPT review file was not available in the named local input locations; the publication order records its operative acceptance and the point-versus-region qualification. We do not invent that missing review or claim GPT inspected this final PDF.

The companion claim ledger maps manuscript labels to accepted claims R01--R23 and immutable source locators. Document build, fidelity and every-page inspection records bind this edition's final PDF and source hashes. They are document/publication checks, not new scientific experiments or symbolic residual tests. The investigation used written proofs and file checks only. No native trajectory, C1 solve, parameter fit, geometry constructor or observer update is needed to rebuild this paper. The portable build helper compiles LaTeX from an existing cache in untrusted mode; it imports no project model. Research licensing follows the repository's existing scope, without a new grant inferred from public visibility.

\begin{thebibliography}{99}\small\raggedright
\bibitem{assessment} Tri-Octagon research programme.
\emph{Native toroidal recirculation and possible diversion: proof review}.
Accepted analytical investigation NATIVE\_TOROIDAL\_RECIRCULATION\_AND\_DIVERSION\_01, 2 November 2026.
\href{evidence/accepted_proof.md}{Retained original proof} and
\href{evidence/accepted_claim_source_ledger.md}{original claim/source ledger}.
The \href{evidence/publication_authorization.md}{publication authorization} records GPT's underlying analytical acceptance and the central-region qualification; the separate GPT review file was unavailable locally.

\bibitem{intensity} H. F. Halld\'orsson.
\emph{Native Intensity Accounting and Conditional Carrier Crossings in the Tri-Octagon Model}.
Repository research edition v0.1, 2 November 2026.
\href{https://github.com/pzychozen/trioctagon-physics/tree/a05c545644417b10f1328e203d421349f7b20bf6/papers/NATIVE_INTENSITY_AND_CONDITIONAL_CARRIER_CROSSINGS/publication/v0.1}{Immutable edition}.
Native intensity, potential scope, central boundaries, host-patch dictionary, cut balance and readouts.

\bibitem{pairedII} H. F. Halld\'orsson.
\emph{Paired Driving and Exterior Identifiability on a Conditional Periodic Carrier}.
Paper II, repository research edition v0.1; manuscript dated 19 October 2026.
\href{https://github.com/pzychozen/trioctagon-physics/tree/c881b682ef6a8128dd1dc316d1850b9aeb1e1c8b/SRG_REDISCOVERY/PAPERS/PAIRED_DRIVING_EXTERIOR_IDENTIFIABILITY/publication/v0.1}{Immutable reading edition}.
Foundations, C1 carrier, native sample/hold, unloaded C6 and limitations. The retained manuscript bears its historical draft label.

\bibitem{pairedI} H. F. Halld\'orsson.
\emph{Paired-Gap Geometry and Toroidal Observation Maps for the Folded Tri-Octagon}.
Paper I, repository research edition v0.1.1; manuscript dated 18 October 2026.
\href{https://github.com/pzychozen/trioctagon-physics/tree/c881b682ef6a8128dd1dc316d1850b9aeb1e1c8b/SRG_REDISCOVERY/PAPERS/PAIRED_GAP_GEOMETRY_OBSERVATION/publication/v0.1.1}{Immutable reading edition}.
Finite measurement triangles, normal-crossing circle, projection bound, axis and clearance. The retained manuscript bears its historical draft label.

\bibitem{fixedgap} Tri-Octagon research programme.
\emph{NCM-FIXED-GAP-01 v0.1}. Accepted geometry report, 31 October 2026.
\href{https://github.com/pzychozen/trioctagon-physics/blob/c881b682ef6a8128dd1dc316d1850b9aeb1e1c8b/papers/PHASE_RECONSTRUCTION_AND_STATIONARY_SUPPORT_AXES/publication/v0.1/provenance/primary/fixed_gap_report.md}{Retained report}, with its accepted review in the same primary-source directory.
Central broken arcs and centered-rotation obstruction.

\bibitem{phase} H. F. Halld\'orsson.
\emph{Phase Reconstruction and Stationary Support Axes in the Tri-Octagon Map}.
Repository research edition v0.1, 2 November 2026.
\href{https://github.com/pzychozen/trioctagon-physics/tree/c881b682ef6a8128dd1dc316d1850b9aeb1e1c8b/papers/PHASE_RECONSTRUCTION_AND_STATIONARY_SUPPORT_AXES/publication/v0.1}{Immutable edition}.
Reconstruction, spatial decoder, stationary axes and observer distinctions.
\end{thebibliography}

\end{document}
